\chapter{Zmienne losowe i rozkłady prawdopodobieństwa}
\label{ch:random-variables}

\providecommand{\atlasref}[2]{#2}

\section{Zmienna losowa i jej rozkład}
Przestrzeń wyników może być bardzo złożona: pojedynczym wynikiem
może być pięć kart, pełna trajektoria urządzenia albo nieskończony
ciąg rzutów monetą. W praktyce interesuje nas często tylko jedna
liczba, na przykład liczba czerwonych kart lub czas oczekiwania.
Zmienna losowa jest więc funkcją, która tłumaczy wyniki doświadczenia
na wartości liczbowe. Jej rozkład opisuje, z jakimi
prawdopodobieństwami pojawiają się zbiory tych wartości.

Mierzalność zapewnia, że zdanie „$X$ przyjmuje wartość należącą
do borelowskiego zbioru $B$” jest dopuszczalnym zdarzeniem.
Nie jest to warunek ciągłości. Funkcja $X$ jest ustalona,
a rozkład jej wartości wynika z wyboru funkcji $X$ oraz
miary probabilistycznej na jej dziedzinie.

\begin{definicja}
Niech $(\Omega,\mathcal F,\Prob)$ będzie przestrzenią probabilistyczną.
Zmienną losową rzeczywistą nazywamy funkcję
$X:\Omega\to\R$, dla której $X^{-1}(B)\in\mathcal F$
dla każdego $B\in\mathcal B(\R)$. Jej rozkładem nazywamy miarę
probabilistyczną $\mu_X$ na $\mathcal B(\R)$, zadaną wzorem
$\mu_X(B):=\Prob(X\in B)=\Prob(X^{-1}(B))$.
Stosujemy oznaczenie $\Law(X)=\mu_X$ oraz $X\sim\mu_X$.
Zapis $\{X\in B\}$ oznacza zbiór $\{\omega\in\Omega:X(\omega)\in B\}$.
Miara $\mu_X$ jest obrazem miary $\Prob$ przez $X$
(\atlasref{ME-02::def:pushforward}{miara obrazowa w teorii miary i całki}):
przeciwobrazy zachowują rozłączne sumy przeliczalne, a
$X^{-1}(\R)=\Omega$, więc $\mu_X$ jest przeliczalnie addytywna
i $\mu_X(\R)=1$. Ponadto $\mu_X(\varnothing)=0$,
gdyż $X^{-1}(\varnothing)=\varnothing$.
\end{definicja}

\begin{stwierdzenie}[Kryterium mierzalności]
Niech $(\Omega,\mathcal F,\Prob)$ będzie przestrzenią probabilistyczną.
Funkcja $X:\Omega\to\R$ jest zmienną losową wtedy i tylko wtedy,
gdy $\{X\leq x\}\in\mathcal F$ dla każdego $x\in\R$.
\end{stwierdzenie}
\begin{proof}
Konieczność wynika z borelowskości przedziału $(-\infty,x]$.
Dla dostateczności zauważamy, że przedziały $(-\infty,x]$ generują
$\mathcal B(\R)$
(\atlasref{ME-01::thm:borel-generators-R}{generatory $\sigma$-ciała borelowskiego}).
Rodzina $\{B\subseteq\R:X^{-1}(B)\in\mathcal F\}$
jest $\sigma$-ciałem, gdyż zawiera $\R$ i przeciwobrazy
zachowują dopełnienia i sumy przeliczalne. Z założenia zawiera
każdą półprostą $(-\infty,x]$, a więc zawiera całe $\mathcal B(\R)$.
\end{proof}

\begin{przyklad}
W modelu dwóch kostek z pełnym $\sigma$-ciałem
(rozdział~\ref{ch:conditioning}) funkcja $X(i,j)=i+j$ jest
mierzalna, gdyż każdy przeciwobraz jest podzbiorem przestrzeni
wyników. Jej rozkład jest skupiony na $\{2,\ldots,12\}$,
ponieważ wszystkie wartości $X$ należą do tego zbioru.
Na $([0,1],\mathcal B([0,1]),\lambda_1|_{\mathcal B([0,1])})$,
gdzie ograniczenie miary ma masę $\lambda_1([0,1])=1$,
zmienna $U(\omega)=\omega$ jest ciągła, zatem borelowsko mierzalna,
a $\mu_U(B)=\lambda_1(B\cap[0,1])$ dla $B\in\mathcal B(\R)$.
Jest to rozkład jednostajny na $[0,1]$. Samo $\Omega$ nie przesądza więc
o postaci rozkładu: ważna jest również zmienna losowa.
\end{przyklad}

\begin{definicja}
Niech $\mu$ będzie borelowską miarą probabilistyczną na $\R$.
Nośnikiem topologicznym $\mu$ nazywamy zbiór
$\supp\mu:=\{x\in\R:\mu((x-\varepsilon,x+\varepsilon))>0
\text{ dla każdego }\varepsilon>0\}$.
\end{definicja}

Nośnik jest najmniejszym domkniętym zbiorem o pełnej mierze
(por. \atlasref{ME-20::thm:support-fullmass-20}{twierdzenie o nośniku miary}).
Istotnie, dopełnienie nośnika jest otwarte, gdyż każdy jego
punkt należy do pewnego otwartego przedziału o mierze zero,
a wszystkie punkty tego przedziału leżą poza nośnikiem.
Wybierając wewnątrz takich przedziałów przedziały o końcach
w $\mathbb Q$, przedstawiamy dopełnienie nośnika jako
co najwyżej przeliczalną sumę przedziałów o mierze zero.
Stąd $\mu(\supp\mu)=1$. Jeżeli domknięty zbiór $C$ ma
pełną miarę, to każdy punkt spoza $C$ ma otwarte otoczenie
o mierze zero, więc $\supp\mu\subseteq C$.

\begin{uwaga}
Rozkład $\mu_X$ zawiera pełną informację o prawdopodobieństwach
zdarzeń wyrażalnych przez wartości $X$, ale nie o zależności $X$
od innych zmiennych. Na przykład na przestrzeni dwóch niezależnych rzutów uczciwą
monetą, kodowanych przez $0$ i $1$, niech $X$ będzie pierwszą
współrzędną, $Y=X$, a $Z$ drugą współrzędną. Wszystkie trzy
zmienne mają ten sam rozkład, lecz $X$ i $Z$ są niezależne,
podczas gdy $X$ i $Y$ nie: $\Prob(X=0,Y=0)=1/2\neq1/4$.
\end{uwaga}

\section{Dystrybuanta}
Rozkład można opisywać przez prawdopodobieństwo każdego zbioru
borelowskiego, lecz w wymiarze jednym wygodniej zebrać
informację w jednej funkcji $F(x)=\Prob(X\leq x)$.
Przechowuje ona prawdopodobieństwa wszystkich lewych półprostych,
a z ich różnic odzyskujemy masy przedziałów. Skok dystrybuanty
w punkcie oznacza dodatnie prawdopodobieństwo dokładnie tej wartości.
Prawostronna ciągłość nie jest dodatkową umową: wynika z
przeliczalnej addytywności miary.

\begin{definicja}
Dystrybuantą rzeczywistej zmiennej losowej $X$ nazywamy funkcję
$F_X:\R\to[0,1]$ określoną wzorem
$F_X(x):=\Prob(X\leq x)=\mu_X((\! -\infty,x])$.
\end{definicja}

\begin{twierdzenie}[Własności dystrybuanty]
Niech $X$ będzie rzeczywistą zmienną losową o dystrybuancie $F$.
Wówczas $F$ jest niemalejąca, prawostronnie ciągła oraz
$\lim_{x\to-\infty}F(x)=0$ i $\lim_{x\to+\infty}F(x)=1$.
Dla dowolnych $a<b$ zachodzą ponadto
\[
\Prob(a<X\leq b)=F(b)-F(a),\qquad
\Prob(X=a)=F(a)-F(a-),
\]
gdzie $F(a-):=\lim_{x\uparrow a}F(x)$.
\end{twierdzenie}
\begin{proof}
Jeśli $x\leq y$, to każdy wynik spełniający $X\leq x$
spełnia też $X\leq y$. Zatem
$\{X\leq x\}\subseteq\{X\leq y\}$ i monotoniczność
miary daje $F(x)\leq F(y)$.

Ustalmy $x\in\R$ i weźmy dowolny ciąg $x_n\downarrow x$.
Zdarzenia $A_n=\{X\leq x_n\}$ maleją. Ich przecięcie
to $\{X\leq x\}$: jeśli $X(\omega)>x$, to dla
dostatecznie dużego $n$ mamy $x_n<X(\omega)$,
więc $\omega\notin A_n$. Odwrotna inkluzja wynika z
$x\leq x_n$. Z ciągłości miary z góry (miara jest skończona)
wynika $F(x_n)\to F(x)$. Jest to prawostronna ciągłość.

Dla $x_n=-n$ zbiory $\{X\leq-n\}$ maleją
do zbioru pustego, ponieważ $X$ przyjmuje tylko
wartości rzeczywiste. Ich prawdopodobieństwa dążą do zera.
Podobnie zbiory $\{X\leq n\}$ rosną do $\Omega$,
więc ich prawdopodobieństwa dążą do jedności.
Monotoniczność $F$ pozwala otrzymane granice dla
ciągów $\pm n$ przenieść na całe półproste.

Dla $a<b$ mamy rozłączny rozkład
$\{X\leq b\}=\{X\leq a\}\mathbin{\dot\cup}
\{a<X\leq b\}$, więc
$\Prob(a<X\leq b)=F(b)-F(a)$.
Wreszcie weźmy $x_n=a-1/n$, $n\geq1$. Wtedy $x_n<a$
i $x_n\uparrow a$. Zachodzi
$\{X\leq x_n\}\uparrow\{X<a\}$: każda wartość mniejsza od $a$
jest nie większa od pewnego $x_n$, a żadna wartość równa $a$
nie należy do tych zdarzeń.
Z monotoniczności i ograniczoności $F$ istnieje granica
lewostronna $F(a-)=\sup_{t<a}F(t)$. Dla każdego $t<a$
mamy $x_n\geq t$ dla dostatecznie dużego $n$, a więc
$F(t)\leq\lim_n F(x_n)\leq F(a-)$. Biorąc supremum po $t<a$,
otrzymujemy $F(x_n)\to F(a-)$. Z ciągłości miary z dołu dostajemy
$F(a-)=\Prob(X<a)$. Ponieważ
$\{X\leq a\}=\{X<a\}\mathbin{\dot\cup}\{X=a\}$,
otrzymujemy $\Prob(X=a)=F(a)-F(a-)$.
\end{proof}

\begin{wniosek}
Dystrybuanta jest ciągła w punkcie $a$ wtedy i tylko wtedy,
gdy $\Prob(X=a)=0$. Wynika to z prawostronnej ciągłości oraz wzoru na skok.
W przeciwieństwie do gęstości dystrybuanta nie jest określona
jedynie prawie wszędzie: dwie prawostronnie ciągłe funkcje
równe $\lambda_1$-prawie wszędzie są równe wszędzie.
Istotnie, zbiór ich równości jest gęsty, a dla dowolnego $a$
można wybrać z niego ciąg punktów większych od $a$ zbieżny do $a$,
i zastosować prawostronną ciągłość.
\end{wniosek}

\begin{twierdzenie}[Charakteryzacja dystrybuant]
Funkcja $F:\R\to\R$ jest dystrybuantą pewnej rzeczywistej
zmiennej losowej wtedy i tylko wtedy, gdy jest niemalejąca,
prawostronnie ciągła, ma granicę $0$ przy $-\infty$
i granicę $1$ przy $+\infty$.
\end{twierdzenie}
\begin{proof}
Konieczność wszystkich warunków wynika z poprzedniego
twierdzenia. Pokażemy, że są one także wystarczające,
konstruując zmienną o zadanej dystrybuancie.
Monotoniczność i obie granice zapewniają $0\leq F(x)\leq1$
dla każdego $x\in\R$.
Niech $\Omega=(0,1)$, $\mathcal F=\mathcal B((0,1))$,
a $\Prob=\lambda_1|_{\mathcal F}$, co daje przestrzeń
probabilistyczną, ponieważ $\lambda_1((0,1))=1$.
Dla $u\in(0,1)$ zdefiniujmy
$Q(u)=\inf\{x\in\R:F(x)\geq u\}$.
Z granic $F$ wynika, że dla każdego $u\in(0,1)$
istnieją $a,b\in\R$ z $F(a)<u\leq F(b)$.
Zbiór pod infimum jest więc niepusty i zawiera się
w $(a,\infty)$; liczba $Q(u)$ jest skończona.

Sprawdźmy kluczową równoważność. Jeżeli $u\leq F(x)$,
to $x$ należy do zbioru w definicji i $Q(u)\leq x$.
Odwrotnie, jeśli $Q(u)\leq x$, to z definicji infimum
dla każdego $n$ istnieje $t_n$ spełniające
$Q(u)\leq t_n<Q(u)+1/n$ oraz $F(t_n)\geq u$.
Z monotoniczności $F$ mamy $F(Q(u)+1/n)\geq u$,
więc możemy wziąć $x_n=Q(u)+1/n$. Prawostronna ciągłość daje
$F(Q(u))\geq u$, a monotoniczność implikuje
$F(x)\geq F(Q(u))\geq u$.
Otrzymaliśmy zatem dla wszystkich $x\in\R$, $u\in(0,1)$
\[
Q(u)\leq x\quad\Longleftrightarrow\quad u\leq F(x).
\]
Ponieważ dla każdego $x$ zachodzi
$\{u\in(0,1):Q(u)\leq x\}=(0,1)\cap(-\infty,F(x)]$,
funkcja $Q$ jest mierzalna na mocy kryterium mierzalności.
Jest również niemalejąca: dla $u<v$ zbiór $\{F\geq v\}$
jest podzbiorem $\{F\geq u\}$, więc $Q(u)\leq Q(v)$.
Dla $U(\omega)=\omega$ mamy
$\Prob(Q(U)\leq x)=\Prob(U\leq F(x))=F(x)$,
ponieważ $U$ jest jednostajna na $(0,1)$, a $0\leq F(x)\leq1$.
Zmienna $X=Q(U)$ realizuje więc żądaną dystrybuantę.
\end{proof}

Dystrybuanta wyznacza rozkład jednoznacznie: półproste
$(-\infty,x]$, $x\in\R$, tworzą układ $\pi$ generujący
$\mathcal B(\R)$, a dwie miary probabilistyczne o tej samej
dystrybuancie zgadzają się na tym układzie i mają tę samą
masę całkowitą. Stosujemy
\atlasref{ME-02::thm:uniqueness-finite}{twierdzenie o jednoznaczności miar skończonych}.

\begin{definicja}
Dla zmiennej losowej $X$ o dystrybuancie $F_X$ i $u\in(0,1)$
kwantylem rzędu $u$ nazywamy
$Q_X(u):=\inf\{x\in\R:F_X(x)\geq u\}$.
Medianą nazywamy dowolną liczbę $m$ spełniającą
$\Prob(X\leq m)\geq1/2$ i $\Prob(X\geq m)\geq1/2$.
\end{definicja}

\begin{uwaga}
Kwantyl $q=Q_X(u)$ spełnia $F_X(q)\geq u$ oraz $F_X(q-)\leq u$.
Pierwsza nierówność wynika z dowodu charakteryzacji dystrybuant;
dla $x<q$ mamy $F_X(x)<u$, więc przejście do granicy z lewej
daje drugą. W szczególności dla $u=1/2$ zachodzi
$\Prob(X\leq q)\geq1/2$ oraz
$\Prob(X\geq q)=1-F_X(q-)\geq1/2$.
Kwantyl $Q_X(1/2)$ jest więc medianą, ale mediana nie musi być jedyna.
Przykładowo zmienna przyjmująca $0$ i $2$ z prawdopodobieństwami
po $1/2$ ma za zbiór median dokładnie $[0,2]$:
dla $m$ z tego przedziału obie wymagane masy są co najmniej
$1/2$, dla $m<0$ pierwsza masa jest zerowa, a dla $m>2$
druga masa jest zerowa.
\end{uwaga}

\section{Rozkłady dyskretne, bezwzględnie ciągłe i osobliwe}
W przypadku dyskretnym przypisujemy masy poszczególnym punktom.
W przypadku rozkładu mającego gęstość zamiast mas punktowych
opisujemy prawdopodobieństwa \emph{przedziałów} przez całkę.
Wartość $f(x)$ nie jest prawdopodobieństwem, że $X=x$:
może przekraczać jeden, podczas gdy $\Prob(X=x)=0$.
Gęstość ma znaczenie lokalne: w punkcie ciągłości skończonej
wersji $f$ zachodzi $\Prob(x<X\leq x+h)=hf(x)+o(h)$
dla $h\downarrow0$, ponieważ
\[
\left|\frac{1}{h}\int_{(x,x+h]}f(t)\,\dd\lambda_1(t)-f(x)\right|
\leq\sup_{x<t\leq x+h}|f(t)-f(x)|\longrightarrow0.
\]
Nie każdy rozkład o ciągłej dystrybuancie posiada gęstość.
Rozkład Cantora
wykaże, dlaczego to rozróżnienie jest ważne.

\begin{definicja}
Rozkład $\mu$ na $\R$ nazywamy dyskretnym, jeśli istnieje zbiór
co najwyżej przeliczalny $D\subseteq\R$ taki, że $\mu(D)=1$.
Dla $x\in\R$ piszemy wtedy $p(x):=\mu(\{x\})$, a funkcję
$p:\R\to[0,1]$ nazywamy funkcją masy prawdopodobieństwa.
\end{definicja}

Jeżeli $\mu$ jest dyskretna, niech $D_0$ będzie co najwyżej
przeliczalnym zbiorem o pełnej mierze. Zbiór
$D=\{x:\mu(\{x\})>0\}$ zawiera się w $D_0$, ponieważ
$\mu(\R\setminus D_0)=0$. Ponadto $D_0\setminus D$
jest przeliczalną sumą zbiorów jednoelementowych o mierze zero.
Stąd $\mu(D)=1$. Przeliczalna addytywność daje
\[
\sum_{x\in D}p(x)=1,\qquad
\mu(B)=\sum_{x\in B\cap D}p(x)\quad(B\in\mathcal B(\R)).
\]
Odwrotnie, dla co najwyżej przeliczalnego $D$ i nieujemnych
mas $p(x)$ o sumie jeden powyższy wzór definiuje rozkład
dyskretny, jako miarę ważoną punktową
(\atlasref{ME-02::ex:weighted-discrete}{teoria miary i całki}).

\begin{definicja}
Borelowską miarę probabilistyczną $\mu$ na $\R$ nazywamy
bezwzględnie ciągłą względem miary Lebesgue'a $\lambda_1$
(zapis $\mu\ll\lambda_1$), gdy $\lambda_1(B)=0$ implikuje $\mu(B)=0$ dla każdego
$B\in\mathcal B(\R)$. Miara $\mu$ jest skończona, a
$\lambda_1$ jest $\sigma$-skończona na tym samym $\sigma$-ciele.
Na mocy \atlasref{ME-18::thm:RN-sigfin-18}{twierdzenia Radona--Nikodyma}
istnieje borelowsko mierzalna funkcja $f:\R\to[0,\infty]$ taka, że
\begin{equation}\label{eq:density}
\mu(B)=\int_B f(x)\,\dd\lambda_1(x)
\quad(B\in\mathcal B(\R)),\qquad
\int_\R f(x)\,\dd\lambda_1(x)=1.
\end{equation}
Funkcję $f$ nazywamy gęstością prawdopodobieństwa; jest ona
wyznaczona jednoznacznie z dokładnością do równości
$\lambda_1$-prawie wszędzie. Możemy przyjąć wersję skończoną
wszędzie, zastępując wartość nieskończoną przez zero na zbiorze
miary Lebesgue'a zero. Odwrotnie, każda nieujemna borelowsko
mierzalna funkcja o całce jeden definiuje przez \eqref{eq:density}
bezwzględnie ciągłą miarę probabilistyczną
(\atlasref{ME-17::ex:ac-density-17}{miara określona gęstością}).
\end{definicja}

\begin{stwierdzenie}
Jeżeli $X$ ma gęstość $f$, to
$F_X(x)=\int_{(-\infty,x]} f(t)\,\dd\lambda_1(t)$.
Dystrybuanta $F_X$ jest bezwzględnie ciągła na $\R$,
a $F_X'(x)=f(x)$
dla $\lambda_1$-prawie każdego $x$.
\end{stwierdzenie}
\begin{proof}
Z definicji gęstości otrzymujemy bezpośrednio
\[
F_X(x)=\mu_X((\!-\infty,x])
=\int_{(-\infty,x]}f(t)\,\dd\lambda_1(t).
\]
Pojedynczy punkt $\{x\}$ ma miarę Lebesgue'a zero,
więc z bezwzględnej ciągłości rozkładu
$\Prob(X=x)=\mu_X(\{x\})=0$.
Z twierdzenia o skokach dystrybuanty wiemy,
że $F_X(x)-F_X(x-)=\Prob(X=x)=0$.
Ponieważ dystrybuanta jest zawsze prawostronnie ciągła,
nie ma również nieciągłości z lewej strony; jest więc ciągła.

Funkcja $f$ jest całkowalna na $\R$, a więc lokalnie
całkowalna. Dla każdego $a\in\R$ możemy zapisać
$F_X(x)=F_X(a)+\int_a^x f(t)\,\dd\lambda_1(t)$,
przyjmując odpowiedni znak całki dla $x<a$.
Stosując na każdym przedziale $[-n,n]$, $n\geq1$,
\atlasref{ME-25::thm:FTC-lebesgue-25}{podstawowe twierdzenie rachunku całkowego dla całki Lebesgue'a},
otrzymujemy $F_X'(x)=f(x)$ poza zbiorem miary zero
w $(-n,n)$. Przeliczalna suma tych zbiorów wyjątkowych ma
miarę zero, więc równość zachodzi prawie wszędzie na $\R$.
Aby uzyskać bezwzględną ciągłość na całej prostej, weźmy
dowolną skończoną rodzinę parami rozłącznych przedziałów
otwartych $(a_j,b_j)$. Wtedy
\[
\sum_j|F_X(b_j)-F_X(a_j)|
=\int_{\bigcup_j(a_j,b_j)} f(t)\,\dd\lambda_1(t).
\]
Punkty końcowe nie zmieniają całki. Z
\atlasref{ME-17::cor:uniform-abs-cont-int-17}{bezwzględnej ciągłości całki względem zbioru całkowania}
dla każdego $\varepsilon>0$ istnieje $\delta>0$ takie,
że prawa strona jest mniejsza od $\varepsilon$, gdy
$\sum_j(b_j-a_j)<\delta$. Jest to bezwzględna ciągłość $F_X$.
Odwrotnie, jeśli dystrybuanta jest bezwzględnie ciągła
na każdym zwartym przedziale, to rozkład ma gęstość równą
jej pochodnej prawie wszędzie; jest to
\atlasref{ME-25::thm:cdf-derivative-25}{twierdzenie o dystrybuancie i bezwzględnej ciągłości miary}.
Samo istnienie pochodnej prawie wszędzie nie wystarcza:
w przypadku rozkładu osobliwego pochodna jest zerowa
prawie wszędzie, co również wynika z przywołanego twierdzenia.
\end{proof}

\begin{uwaga}
Sama ciągłość dystrybuanty nie oznacza istnienia gęstości.
Istnieją \emph{rozkłady osobliwe bezatomowe}, na przykład rozkład
Cantora: jest skupiony na zbiorze Cantora o mierze Lebesgue'a zero,
a jednocześnie nie ma atomów. Jego dystrybuanta jest ciągła,
ale nie jest bezwzględnie ciągła. Ogólna miara probabilistyczna na $\R$
może zawierać jednocześnie część dyskretną, bezwzględnie ciągłą
i osobliwą bezatomową.
\end{uwaga}


\begin{twierdzenie}[Rozkład miary na część dyskretną, bezwzględnie ciągłą i osobliwą bezatomową]
Niech $\mu$ będzie borelowską miarą probabilistyczną na $\R$.
Wówczas istnieją jednoznacznie wyznaczone nieujemne
borelowskie miary skończone
$\mu_{\mathrm d},\mu_{\mathrm{ac}},\mu_{\mathrm{sc}}$
takie, że $\mu=\mu_{\mathrm d}+\mu_{\mathrm{ac}}+
\mu_{\mathrm{sc}}$, gdzie $\mu_{\mathrm d}$ jest skupiona
na zbiorze co najwyżej przeliczalnym,
$\mu_{\mathrm{ac}}\ll\lambda_1$,
a $\mu_{\mathrm{sc}}\perp\lambda_1$ i nie ma atomów.
Zapis $\eta\perp\lambda_1$ oznacza, że istnieje borelowski
zbiór $S$ taki, że $\lambda_1(S)=0$ i $\eta(\R\setminus S)=0$.
Por. \atlasref{ME-17::thm:atomic-cont-singular-17}{rozkład części osobliwej w teorii miary i całki}.
Każda z trzech miar może mieć masę mniejszą od jedności,
a nawet być miarą zerową.
\end{twierdzenie}
\begin{proof}
Najpierw zbiór atomów punktowych
$D=\{x\in\R:\mu(\{x\})>0\}$ jest co najwyżej
przeliczalny: dla każdego $m\geq1$ istnieje najwyżej
$m$ punktów o masie większej niż $1/m$, a
$D$ jest sumą tych skończonych zbiorów.
Zbiór $D$ jest borelowski jako co najwyżej przeliczalny.
Definiujemy dla $B\in\mathcal B(\R)$ miary ograniczone
$\mu_{\mathrm d}(B)=\mu(B\cap D)$
oraz $\rho(B)=\mu(B\setminus D)$.
Mamy $\mu=\mu_{\mathrm d}+\rho$ oraz $\rho(\{x\})=0$ dla
każdego $x\in\R$. Na prostej skończona miara borelowska
bez dodatnich mas punktowych jest bezatomowa
(\atlasref{ME-17::thm:atomic-cont-singular-17}{twierdzenie o rozkładzie na część atomową i bezatomową}).

Wyprowadzimy teraz potrzebny rozkład na część bezwzględnie
ciągłą i osobliwą z przyjętego twierdzenia
Radona--Nikodyma. Połóżmy $\nu=\rho+\lambda_1$.
Miara $\nu$ jest $\sigma$-skończona, ponieważ
$\rho$ jest skończona, a $\lambda_1$ jest
$\sigma$-skończona (każdy przedział $[-n,n]$ ma skończoną
miarę $\nu$). Ponieważ $\lambda_1\leq\nu$ i $\rho\leq\nu$,
obie miary są bezwzględnie ciągłe względem $\nu$.
Wystarczy użyć borelowskiej pochodnej
$g=\dd\lambda_1/\dd\nu$, której istnienie wynika z twierdzenia
Radona--Nikodyma. Dla $S=\{x:g(x)=0\}$ mamy
$\lambda_1(S)=\int_Sg\,\dd\nu=0$.
Definiujemy dla $B\in\mathcal B(\R)$
\[
\mu_{\mathrm{sc}}(B)=\rho(B\cap S),\qquad
\mu_{\mathrm{ac}}(B)=\rho(B\setminus S).
\]
Są to nieujemne miary skończone, a ich suma jest równa $\rho$.
Pierwsza miara jest osobliwa, bo skupia się na $S$.
Druga jest bezwzględnie ciągła względem $\lambda_1$:
jeśli $\lambda_1(B)=0$, to
$\int_B g\,\dd\nu=0$, a więc $g=0$
$\nu$-prawie wszędzie na $B$. Stąd
$\nu(B\cap\{g>0\})=0$, a z $\rho\leq\nu$ wynika
$\mu_{\mathrm{ac}}(B)=\rho(B\cap\{g>0\})=0$.
Obie części mają zerowe masy punktowe, gdyż są zdominowane
przez $\rho$, a z przywołanego twierdzenia o miarach
na prostej wynika ich bezatomowość.

Dla jednoznaczności zauważmy najpierw, że cała
masa punktowa musi należeć do $\mu_{\mathrm d}$,
gdyż dwie pozostałe części mają zerowe masy punktowe.
Każda dopuszczalna część dyskretna ma więc masę
$\mu(\{x\})$ w punkcie $x$ i jest skupiona na $D$,
zatem dla każdego borelowskiego $B$ ma wartość
$\sum_{x\in B\cap D}\mu(\{x\})$. Wyznacza to ją jednoznacznie.
Przypuśćmy następnie, że
$\rho=\alpha_1+\beta_1=\alpha_2+\beta_2$,
gdzie $\alpha_i\ll\lambda_1$, a $\beta_i$
są skupione na zbiorach $S_i$ o miarze Lebesgue'a zero.
Dla $S=S_1\cup S_2$ obie miary $\beta_i$
zanikają poza $S$, więc $\alpha_1$ i $\alpha_2$
są tam równe $\rho$. Na samym $S$ obie zanikają
z bezwzględnej ciągłości względem $\lambda_1$.
Zatem $\alpha_1=\alpha_2$ oraz
$\beta_1=\beta_2$. To kończy dowód jednoznaczności.
\end{proof}

\begin{przyklad}[Gęstość większa niż jeden i zerowa masa punktu]
Niech $X$ będzie jednostajna na $(0,1/4)$. Jej gęstość
to $f_X(x)=4\chi_{(0,1/4)}(x)$. Mamy $f_X(1/8)=4$,
co nie przeczy aksjomatom, ponieważ gęstość nie jest
prawdopodobieństwem zdarzenia $\{X=1/8\}$.
W rzeczywistości $\Prob(X=1/8)=0$, zaś
$\Prob(0<X<1/8)=\int_0^{1/8}4\,\dd\lambda_1(x)=1/2$.
Jeśli wartości $X$ mają jednostkę fizyczną, gęstość ma
jednostkę odwrotną. Prawdopodobieństwem jest całka gęstości
po zbiorze wartości, nie wartość gęstości w punkcie.
\end{przyklad}

\begin{przyklad}[Rozkład Cantora i brak gęstości]
Na przestrzeni ciągów Bernoulliego z parametrem $1/2$
(rozdział~\ref{ch:probability}) niech $B_n$ oznacza $n$-tą
współrzędną. Zmienne $B_n$ są niezależne, przyjmują wartości
w $\{0,1\}$ i $\Prob(B_n=0)=\Prob(B_n=1)=1/2$.
Każda współrzędna jest mierzalna, ponieważ jej przeciwobrazy
zbiorów borelowskich są cylindrami, zbiorem pustym lub całą
przestrzenią. Dla parami różnych $n_1,\ldots,n_k$ i
$b_1,\ldots,b_k\in\{0,1\}$ konstrukcja miary daje
$\Prob(B_{n_1}=b_1,\ldots,B_{n_k}=b_k)=2^{-k}$.
Dla dowolnych borelowskich $A_1,\ldots,A_k$ sumujemy tę
równość po słowach z $\prod_{j=1}^k(A_j\cap\{0,1\})$.
Dostajemy iloczyn $\prod_{j=1}^k\Prob(B_{n_j}\in A_j)$,
co uzasadnia niezależność zmiennych losowych.
Szereg $X=\sum_{n=1}^\infty 2B_n/3^n$ zbiega bezwzględnie
dla każdego wyniku, gdyż $\sum_{n=1}^\infty 2/3^n=1$.
Zmienna $X$ jest mierzalna jako granica mierzalnych sum częściowych
(\atlasref{ME-06::thm:pointwise-limits}{twierdzenie o granicach funkcji mierzalnych}).
Dla każdego całkowitego $m\geq1$ i ustalonego słowa
$b=(b_1,\ldots,b_m)\in\{0,1\}^m$ połóżmy
\[
a_b=\sum_{n=1}^m\frac{2b_n}{3^n},\qquad
I_b=[a_b,a_b+3^{-m}].
\]
Na zdarzeniu $\{B_1=b_1,\ldots,B_m=b_m\}$ mamy $X\in I_b$,
gdyż ogon szeregu spełnia
$0\leq\sum_{n=m+1}^\infty 2B_n/3^n\leq
\sum_{n=m+1}^\infty 2/3^n=3^{-m}$.
Przedziały $I_b$ są dokładnie przedziałami $m$-tego etapu
konstrukcji zbioru Cantora: przejście do następnego etapu
zastępuje każdy przedział jego lewą i prawą trzecią częścią.
Na każdym ustalonym etapie przedziały są parami rozłączne,
oddzielone przerwami o dodatniej długości.
Ich suma $C_m$ składa się z $2^m$ domkniętych przedziałów
długości $3^{-m}$, stąd $\lambda_1(C_m)=2^m3^{-m}=(2/3)^m$.
Zbiory $C_m$ maleją, są borelowskie i $\lambda_1(C_1)<\infty$.
Zatem $C=\bigcap_{m\geq1}C_m$ zawiera wszystkie wartości $X$,
a $\lambda_1(C)=0$ z ciągłości miary z góry.
Mamy $\Prob(X\in C)=1$, więc rozkład $X$ nie może mieć
gęstości względem $\lambda_1$.

Jeśli $x\in C$, to na każdym etapie należy do dokładnie
jednego $I_b$. Zdarzenie $\{X=x\}$ jest podzbiorem
$\{B_1=b_1,\ldots,B_m=b_m\}$, którego prawdopodobieństwo
z niezależności wynosi $2^{-m}$. Stąd
$\Prob(X=x)\leq2^{-m}$ dla każdego $m$ i w granicy
$\Prob(X=x)=0$. Poza $C$ również wszystkie masy punktowe
są zerowe.
Z twierdzenia o skokach dystrybuanty wynika, że $F_X$ jest
ciągła. Rozkład $X$ jest osobliwy względem miary Lebesgue'a
i bezatomowy na mocy przywołanego wcześniej twierdzenia
o skończonych miarach borelowskich na prostej.
\end{przyklad}

\begin{przyklad}[Rozkład mieszany]
Niech $B$ przyjmuje wartości $0$ i $1$ z prawdopodobieństwami po $1/2$,
a $U$ jest od niego niezależna i jednostajna na $(0,1)$.
Definiujemy $X=BU$, co jest zmienną losową jako iloczyn
dwóch rzeczywistych funkcji mierzalnych.
Zdarzenie $\{U\in(0,1)\}$ ma pełną
miarę, więc $\Prob(X=0)=\Prob(B=0)=1/2$.
Dla $0<x<1$ z rozkładu na zdarzenia $\{B=0\}$ i $\{B=1\}$
oraz niezależności otrzymujemy
\[
F_X(x)=\Prob(B=0)+\Prob(B=1,U\leq x)=\frac12+\frac{x}{2}.
\] Łącznie
\[
F_X(x)=\begin{cases}
0,&x<0,\\
\frac12,&x=0,\\
\frac{1+x}{2},&0<x<1,\\
1,&x\geq1.
\end{cases}
\]
Rozkład nie ma gęstości względem $\lambda_1$, gdyż ma atom w zerze,
ale jego część bezwzględnie ciągła ma gęstość
$\frac12\chi_{(0,1)}$. Dokładniej, dla borelowskiego $A$
z niezależności i rozbicia według $B$ mamy
$\mu_X(A)=\frac12\delta_0(A)+\frac12\lambda_1(A\cap(0,1))$,
gdzie $\delta_0$ jest miarą Diraca w zerze.
\end{przyklad}

\section{Wektory losowe i rozkłady łączne}
Znajomość rozkładu wzrostu i rozkładu masy ciała osobno nie
mówi, czy wzrost i masa są ze sobą związane. Potrzebny jest
\emph{rozkład łączny} wektora dwóch cech. Rozkłady brzegowe
powstają przez rzutowanie miary łącznej na wybrane współrzędne.
Jeżeli istnieje gęstość łączna, rzutowaniu odpowiada
całkowanie po pomijanych współrzędnych. Odwrotna
rekonstrukcja rozkładu łącznego z samych rozkładów brzegowych
na ogół nie jest możliwa; niezależność stanowi szczególny
przypadek, gdy taka rekonstrukcja jest możliwa przez iloczyn.

\begin{definicja}
Niech $(\Omega,\mathcal F,\Prob)$ będzie przestrzenią probabilistyczną,
a $d\geq1$ liczbą całkowitą. Wektorem losowym w $\R^d$ nazywamy funkcję
$X=(X_1,\ldots,X_d):\Omega\to\R^d$ mierzalną względem
$\mathcal F$ i $\mathcal B(\R^d)$. Jej rozkładem łącznym jest
$\mu_X(B)=\Prob(X\in B)$ dla $B\in\mathcal B(\R^d)$.
Jest to miara probabilistyczna jako obraz miary $\Prob$
przez mierzalne odwzorowanie $X$.
Rozkład $X_i$ nazywamy $i$-tym rozkładem brzegowym.
Jeżeli $\pi_i:\R^d\to\R$ jest rzutem na $i$-tą współrzędną,
to dla każdego borelowskiego $A\subseteq\R$
\[
\mu_{X_i}(A)=\Prob(X_i\in A)
=\Prob(X\in\pi_i^{-1}(A))=\mu_X(\pi_i^{-1}(A)).
\]
Zatem $\mu_{X_i}$ jest miarą obrazową $\mu_X$ przez $\pi_i$.
\end{definicja}

\begin{stwierdzenie}
Niech $X=(X_1,\ldots,X_d):\Omega\to\R^d$ będzie funkcją
na przestrzeni probabilistycznej $(\Omega,\mathcal F,\Prob)$.
Wektor $X$ jest mierzalny wtedy i tylko wtedy,
gdy każda współrzędna $X_i$ jest rzeczywistą zmienną losową.
\end{stwierdzenie}
\begin{proof}
Jeżeli $X$ jest mierzalny, to $X_i=\pi_i\circ X$ jest mierzalna
jako złożenie funkcji mierzalnych: rzut $\pi_i$ jest ciągły,
a zatem borelowsko mierzalny.
W przeciwnym kierunku przeciwobraz prostokąta
$B_1\times\cdots\times B_d$, gdzie $B_i\in\mathcal B(\R)$,
jest mierzalnym przecięciem $\bigcap_{i=1}^d\{X_i\in B_i\}$.
Rodzina zbiorów o mierzalnym przeciwobrazie przez $X$ jest
$\sigma$-ciałem, gdyż zawiera $\R^d$ (jego przeciwobraz
to $\Omega$), a przeciwobrazy zachowują dopełnienia
i sumy przeliczalne. Zawiera ona wszystkie borelowskie
prostokąty, które generują $\mathcal B(\R^d)$
(\atlasref{ME-01::thm:borel-Rd}{generatory $\sigma$-ciała borelowskiego w przestrzeni euklidesowej}).
Zatem $X$ jest mierzalny.
\end{proof}

\begin{definicja}
Rzeczywiste zmienne losowe $X_1,\ldots,X_d$ określone na tej
samej przestrzeni probabilistycznej nazywamy niezależnymi,
jeżeli dla wszystkich $B_1,\ldots,B_d\in\mathcal B(\R)$ zachodzi
$\Prob(X_1\in B_1,\ldots,X_d\in B_d)=\prod_{i=1}^d\Prob(X_i\in B_i)$.
Dla rodziny $(X_i)_{i\in I}$ niezależność oznacza spełnienie tej
własności przez każdą skończoną rodzinę o parami różnych
indeksach $i_1,\ldots,i_m\in I$. Warunek dla $d$ zmiennych
implikuje warunki dla wszystkich ich podrodzin, gdyż dla
pomijanych współrzędnych można wziąć $B_i=\R$.
\end{definicja}

\begin{stwierdzenie}
Niech $X_1,\ldots,X_d$, $d\geq1$, będą rzeczywistymi zmiennymi
losowymi na wspólnej przestrzeni probabilistycznej.
Są one niezależne wtedy i tylko wtedy, gdy ich rozkład łączny
jest miarą produktową
\[
\mu_{(X_1,\ldots,X_d)}=\mu_{X_1}\otimes\cdots\otimes\mu_{X_d}.
\]
\end{stwierdzenie}
\begin{proof}
Miara produktowa istnieje, ponieważ wszystkie czynniki są
miarami probabilistycznymi, zatem skończonymi
(\atlasref{ME-12::thm:product-measure-12}{twierdzenie o mierze produktowej},
stosowane kolejno do skończenie wielu czynników).
Na prostokącie $B_1\times\cdots\times B_d$ jej wartością
jest $\prod_{i=1}^d\mu_{X_i}(B_i)$.
Jeżeli zmienne są niezależne, wartość ta jest równa
$\Prob(X_1\in B_1,\ldots,X_d\in B_d)$.
Prostokąty borelowskie tworzą układ $\pi$ generujący
$\mathcal B(\R^d)$, a obie miary mają masę całkowitą jeden.
Z \atlasref{ME-02::thm:uniqueness-finite}{jednoznaczności miar skończonych}
wynika ich równość. Odwrotnie, z równości miar przez
obliczenie ich wartości na prostokątach otrzymujemy definicję
niezależności.
\end{proof}

\begin{definicja}
Jeśli rozkład wektora losowego $X:\Omega\to\R^d$ jest bezwzględnie
ciągły względem $d$-wymiarowej miary Lebesgue'a $\lambda_d$
na $\mathcal B(\R^d)$, nieujemną borelowsko mierzalną funkcję
$f_X:\R^d\to[0,\infty]$ spełniającą
$\Prob(X\in B)=\int_B f_X(x)\,\dd\lambda_d(x)$ dla każdego
$B\in\mathcal B(\R^d)$ nazywamy gęstością łączną.
Jej istnienie i jednoznaczność $\lambda_d$-prawie wszędzie
wynikają z twierdzenia Radona--Nikodyma, a dla $B=\R^d$
otrzymujemy $\int_{\R^d}f_X(x)\,\dd\lambda_d(x)=1$.
\end{definicja}

\begin{twierdzenie}[Gęstości brzegowe i niezależność]
Niech $(X,Y)$ będzie wektorem losowym o gęstości łącznej
$f_{X,Y}:\R^2\to[0,\infty]$. Wówczas funkcje
\[
f_X(x)=\int_\R f_{X,Y}(x,y)\,\dd\lambda_1(y),\qquad
f_Y(y)=\int_\R f_{X,Y}(x,y)\,\dd\lambda_1(x)
\]
są borelowsko mierzalnymi gęstościami rozkładów brzegowych.
Każdą z nich można zastąpić inną nieujemną funkcją
borelowską równą jej $\lambda_1$-prawie wszędzie,
w szczególności przyjąć wersję wszędzie skończoną.
Na niezupełnym borelowskim $\sigma$-ciele dowolna
modyfikacja na zbiorze miary zero nie musi zachowywać
mierzalności.
W iloczynie gęstości rozszerzonych przyjmujemy $0\cdot\infty=0$. Zmienne $X,Y$ są niezależne wtedy i tylko wtedy,
gdy $f_{X,Y}(x,y)=f_X(x)f_Y(y)$ dla $\lambda_2$-prawie każdego
$(x,y)\in\R^2$.
\end{twierdzenie}
\begin{proof}
Najpierw sprawdźmy, że podane całki są rzeczywiście
gęstościami brzegowymi. Dla borelowskiego $A\subseteq\R$
z definicji rozkładu łącznego
\[
\Prob(X\in A)=\Prob((X,Y)\in A\times\R)
=\int_{A\times\R}f_{X,Y}(x,y)\,\dd\lambda_2(x,y).
\]
Ponieważ funkcja podcałkowa jest nieujemna i borelowsko mierzalna,
a $\lambda_2=\lambda_1\otimes\lambda_1$ na borelowskim
$\sigma$-ciele (\atlasref{ME-12::thm:lebesgue-products-12}{miary Lebesgue'a i produkt}),
możemy zastosować \atlasref{ME-12::thm:tonelli-12}{twierdzenie Tonellego}
i scałkować najpierw po $y$:
\[
\Prob(X\in A)=\int_A
\left(\int_\R f_{X,Y}(x,y)\,\dd\lambda_1(y)\right)
\dd\lambda_1(x).
\]
Twierdzenie Tonellego zapewnia także borelowską mierzalność
funkcji $f_X$. Dla $A=\R$ otrzymujemy
$\int_\R f_X(x)\,\dd\lambda_1(x)=1$, zatem $f_X$ jest
skończona $\lambda_1$-prawie wszędzie. Analogicznie, całkując
najpierw po $x$, uzyskujemy mierzalność, unormowanie i
własność gęstości dla $f_Y$.
Zmiana wersji $f_X$ lub $f_Y$ na zbiorze miary zero
zmienia iloczyn tylko na zbiorze miary $\lambda_2$ zero,
co wynika z twierdzenia Tonellego dla funkcji charakterystycznych
odpowiednich zbiorów produktowych.

Jeśli $f_{X,Y}=f_X f_Y$ prawie wszędzie, to dla
borelowskich $A,B$ twierdzenie Tonellego daje
$\Prob(X\in A,Y\in B)
=(\int_A f_X\,\dd\lambda_1)
 (\int_B f_Y\,\dd\lambda_1)
=\Prob(X\in A)\Prob(Y\in B)$.
Jest to definicja niezależności.

Odwrotnie, przy niezależności z poprzedniego stwierdzenia
$\mu_{(X,Y)}=\mu_X\otimes\mu_Y$.
Miara $\eta(C)=\int_C f_X(x)f_Y(y)\,\dd\lambda_2(x,y)$
jest probabilistyczna: Tonelli daje masę całkowitą $1$.
Dla każdego prostokąta borelowskiego $A\times B$ mamy
\[
\eta(A\times B)=
\left(\int_A f_X(x)\,\dd\lambda_1(x)\right)
\left(\int_B f_Y(y)\,\dd\lambda_1(y)\right)
=\mu_X(A)\mu_Y(B).
\]
Z jednoznaczności miar probabilistycznych na układzie $\pi$
prostokątów dostajemy $\eta=\mu_X\otimes\mu_Y$.
Jednoznaczność pochodnej Radona--Nikodyma daje
$f_{X,Y}=f_Xf_Y$ dla $\lambda_2$-prawie każdego punktu.
\end{proof}

\begin{przyklad}[Jednostajność na trójkącie]
Niech $(X,Y)$ ma gęstość $f_{X,Y}(x,y)=2$ na zbiorze
$T=\{(x,y):0<x<y<1\}$ i zero poza nim.
Z twierdzenia Tonellego
$\lambda_2(T)=\int_0^1(1-x)\,\dd\lambda_1(x)=1/2$,
więc gęstość jest unormowana. Taki wektor istnieje:
na $T$ z borelowskim $\sigma$-ciałem i miarą
$2\lambda_2|_T$ wystarczy wziąć rzuty współrzędne.
Dla $0<x<1$ mamy $f_X(x)=\int_x^1 2\,\dd\lambda_1(y)=2(1-x)$,
a dla $0<y<1$ mamy $f_Y(y)=\int_0^y 2\,\dd\lambda_1(x)=2y$.
Poza $(0,1)$ obie gęstości są zerowe. Nie są stałe na
$(0,1)$, więc rozkłady brzegowe nie są jednostajne na tym przedziale.
Na zbiorze $\{(x,y):0<y<x<1\}$, mającym dodatnią
miarę Lebesgue'a, iloczyn $f_X(x)f_Y(y)=4(1-x)y$ jest
dodatni, lecz $f_{X,Y}(x,y)=0$. Zatem z kryterium iloczynowego
$X,Y$ nie są niezależne. Można to też sprawdzić na zdarzeniach:
$\Prob(X>1/2)=1/4$ oraz $\Prob(Y<1/2)=1/4$
z całkowania gęstości brzegowych, ale
$\Prob(X>1/2,Y<1/2)=0\neq1/16$, ponieważ $\Prob(X<Y)=1$.
\end{przyklad}

\section{Rozkłady warunkowe przy istnieniu gęstości}
Jeżeli rozkład $Y$ jest bezatomowy, to dla każdego $y\in\R$
zdarzenie $\{Y=y\}$ ma prawdopodobieństwo zero.
Wyrażenia $\Prob(X\in A\mid Y=y)$ nie można wówczas
definiować przez iloraz prawdopodobieństw z rozdziału~\ref{ch:conditioning}.
Gdy istnieje gęstość łączna, jej rozkład na czynniki
„gęstość brzegowa razy gęstość warunkowa” stanowi
naturalne rozszerzenie wzoru na warunkowanie. Na tym etapie
konstruujemy je bez wchodzenia w ogólną teorię regularnych
prawdopodobieństw warunkowych.

\begin{stwierdzenie}[Warunkowa gęstość]
Niech $(X,Y)$ będzie wektorem losowym o gęstości łącznej $f_{X,Y}$,
a $f_Y$ niech oznacza gęstość brzegową $Y$.
Wybierzmy nieujemną, wszędzie skończoną borelowską wersję
$f_{X,Y}$. Połóżmy
\[
m(y)=\int_\R f_{X,Y}(x,y)\,\dd\lambda_1(x),\qquad
D=\{y\in\R:0<m(y)<\infty\}.
\]
Jako wersję gęstości brzegowej przyjmijmy $f_Y(y)=m(y)$
dla $m(y)<\infty$ oraz $f_Y(y)=0$ dla $m(y)=\infty$.
Niech $r$ będzie ustaloną nieujemną, skończoną borelowską
gęstością probabilistyczną, na przykład $r=\chi_{(0,1)}$.
Definiujemy
\[
k(y,x):=f_{X\mid Y=y}(x)=
\begin{cases}
f_{X,Y}(x,y)/m(y),& y\in D,\\
r(x),& y\notin D.
\end{cases}
\]
Funkcja $k$ jest wspólnie borelowsko mierzalna, a dla każdego
$y$ funkcja $x\mapsto k(y,x)$ jest gęstością probabilistyczną.
Ponadto $\Law(Y)(\R\setminus D)=0$.
Dla borelowskich $A,B\subseteq\R$
\begin{equation}\label{eq:joint_cond}
\Prob(X\in A,Y\in B)
=\int_B\left(\int_A f_{X\mid Y=y}(x)\,\dd\lambda_1(x)\right)
 f_Y(y)\,\dd\lambda_1(y),
\end{equation}
przy każdym ustalonym wyborze gęstości $r$ jak wyżej.
\end{stwierdzenie}
\begin{proof}
Twierdzenie Tonellego zapewnia borelowską mierzalność $m$ i
$\int_\R m(y)\,\dd\lambda_1(y)=1$. Zatem
$\lambda_1(\{m=\infty\})=0$, a $f_Y$ jest gęstością
$\Law(Y)$ na mocy twierdzenia o gęstościach brzegowych.
Zbiór $D$ jest borelowski. Na jego dopełnieniu $f_Y=0$,
więc $\Law(Y)(\R\setminus D)=\int_{\R\setminus D}f_Y\,\dd\lambda_1=0$.

Na $D\times\R$ iloraz definiujący $k$ jest borelowski
jako iloraz funkcji borelowskich o dodatnim skończonym
mianowniku. Na $(\R\setminus D)\times\R$ funkcja $k(y,x)=r(x)$
jest również borelowska. Oba zbiory są borelowskie, co
dowodzi wspólnej mierzalności. Dla $y\in D$ mamy
$\int_\R k(y,x)\,\dd\lambda_1(x)=m(y)/m(y)=1$,
a dla $y\notin D$ całka jest równa jeden z wyboru $r$.

Dla każdego ustalonego $y$ funkcja zbioru
$K(y,A):=\int_A k(y,x)\,\dd\lambda_1(x)$,
$A\in\mathcal B(\R)$, jest miarą probabilistyczną:
jest to miara określona nieujemną gęstością o całce jeden.
Dla każdego ustalonego borelowskiego $A$ funkcja
$y\mapsto K(y,A)$ jest borelowska na mocy twierdzenia
Tonellego i przyjmuje wartości w $[0,1]$.
Dla $y\in D$ mamy $f_Y(y)=m(y)$, zatem z definicji ilorazu wynika
$k(y,x)f_Y(y)=f_{X,Y}(x,y)$ dla każdego $x\in\R$,
a po scałkowaniu po $A$ otrzymujemy
$K(y,A)f_Y(y)=\int_A f_{X,Y}(x,y)\,\dd\lambda_1(x)$.
Ta równość zachodzi również dla $m(y)=0$, ponieważ obie
strony są wtedy zerowe; zbiór $\{m=\infty\}$ jest
$\lambda_1$-zerowy. Zatem po scałkowaniu po $B$ i zastosowaniu
twierdzenia Tonellego otrzymujemy \eqref{eq:joint_cond}.
\end{proof}

\begin{uwaga}
Symbol $\Prob(X\in A\mid Y=y)$, rozumiany przez powyższy wzór,
nie jest zwykłym ilorazem $\Prob(X\in A,Y=y)/\Prob(Y=y)$,
gdyż przy istnieniu gęstości rozkładu $Y$ licznik i mianownik
są zerowe dla każdego $y$.
Definiujemy ten symbol jako $K(y,A)$ dla wybranej wersji $k$.
Jeżeli $k_1,k_2$ są dwiema wspólnie borelowskimi gęstościami
warunkowymi spełniającymi \eqref{eq:joint_cond}, to dla każdego
borelowskiego $A$ funkcje $K_1(\cdot,A)$ i $K_2(\cdot,A)$
mają równe całki po wszystkich borelowskich $B$ względem
$\Law(Y)$, stąd są równe $\Law(Y)$-prawie wszędzie:
wystarczy rozpatrzyć zbiory
$\{y:K_1(y,A)-K_2(y,A)\geq1/n\}$ oraz
$\{y:K_2(y,A)-K_1(y,A)\geq1/n\}$ dla $n\geq1$.
Każdy z nich ma miarę zero, ponieważ całka różnicy
po nim jest zerowa.
Stosując tę obserwację do przeliczalnej rodziny półprostych
$(-\infty,q]$, $q\in\mathbb Q$, i sumując zbiory wyjątkowe,
uzyskujemy jeden zbiór $\Law(Y)$-miary zero, poza którym
obie miary warunkowe mają dystrybuanty równe na $\mathbb Q$.
Dla każdego $t\in\R$ wybieramy ciąg liczb wymiernych
$q_n\downarrow t$ i z prawostronnej ciągłości otrzymujemy
równość dystrybuant w $t$. Na mocy jednoznaczności rozkładu
wyznaczonego przez dystrybuantę miary te są równe na
wszystkich zbiorach borelowskich.
Ich gęstości są wtedy równe $\lambda_1$-prawie wszędzie
w zmiennej $x$. W szczególności wartości warunkowania w
ustalonym punkcie $y$ nie są wyznaczone jednoznacznie przez
rozkład łączny, jeśli $\Law(Y)(\{y\})=0$: można zastąpić
$k(y,\cdot)$ w tym jednym punkcie dowolną ustaloną
borelowską gęstością, zachowując mierzalność wspólną
i wszystkie tożsamości całkowe \eqref{eq:joint_cond}.
\end{uwaga}

\begin{przyklad}[Warunkowanie na ustaloną wartość zmiennej ciągłej]
Wróćmy do rozkładu jednostajnego na trójkącie
$T=\{(x,y):0<x<y<1\}$, o gęstości łącznej
$f_{X,Y}(x,y)=2\chi_T(x,y)$.
Wykazaliśmy, że $f_Y(y)=2y$ dla $0<y<1$.
Zatem dla każdej takiej wartości $y$ możemy wybrać
wersję gęstości warunkowej
$f_{X\mid Y=y}(x)=y^{-1}\chi_{(0,y)}(x)$.
Jest ona unormowana, gdyż $\int_0^y y^{-1}\,\dd\lambda_1(x)=1$.
Na przykład przy $0<a<y<1$ otrzymujemy
$\Prob(X\leq a\mid Y=y)=a/y$ w sensie tej wersji
rozkładu warunkowego. Wniosek ma prostą interpretację:
dla tej wersji, przy $Y=y$, rozkład warunkowy $X$ jest
jednostajny na przedziale $(0,y)$. Nie obliczaliśmy tu ilorazu
prawdopodobieństw zdarzeń o zerowej mierze;
posłużyliśmy się rozkładem łącznym i jego faktoryzacją.
\end{przyklad}



\begin{twierdzenie}[Warunkowa wartość oczekiwana obliczana z gęstości]
Niech $(X,Y)$ będzie rzeczywistym wektorem losowym
mającym borelowską gęstość łączną $f_{X,Y}$.
Niech $k(y,x)$ będzie wspólnie borelowską, nieujemną
gęstością warunkową spełniającą \eqref{eq:joint_cond},
unormowaną dla każdego $y$, jak w stwierdzeniu o gęstości
warunkowej.
Jeśli $h:\R\to\R$ jest funkcją borelowską i
$\E|h(X)|<\infty$, to zbiór
\[
N=\left\{y:\int_\R |h(x)|k(y,x)\,\dd\lambda_1(x)=\infty\right\}
\]
jest borelowski i $\Law(Y)(N)=0$. Funkcja
\[
g(y)=
\begin{cases}
\displaystyle\int_\R h(x)k(y,x)\,\dd\lambda_1(x),&y\notin N,\\
0,&y\in N
\end{cases}
\]
jest borelowska, a $g(Y)$ jest całkowalna i spełnia
\[
\E[h(X)\mid\sigma(Y)]=g(Y)\quad\text{p.n.}
\]
\end{twierdzenie}
\begin{proof}
Dla dowolnej nieujemnej funkcji borelowskiej $\varphi$
niech $G_\varphi(y)=\int_\R \varphi(x)k(y,x)\,\dd\lambda_1(x)$.
Z twierdzenia Tonellego jest to borelowska funkcja
o wartościach w $[0,\infty]$. Dla każdego borelowskiego
$B\subseteq\R$ zachodzi
\[
\int_\Omega \varphi(X(\omega))\chi_{Y^{-1}(B)}(\omega)\,
\dd\Prob(\omega)
=\int_B G_\varphi(y)\,\dd\Law(Y)(y).
\]
Dla $\varphi=\chi_A$ lewa strona jest równa
$\Prob(X\in A,Y\in B)$, a prawa
$\int_B K(y,A)f_Y(y)\,\dd\lambda_1(y)$ na mocy
wzoru na całkę względem miary określonej gęstością
(\atlasref{ME-18::lem:integration-density-18}{całkowanie względem miary z gęstością}).
Jest to \eqref{eq:joint_cond}; przypadek nieujemnych
funkcji prostych wynika z liniowości.
Wybierając nieujemne funkcje proste $\varphi_n\uparrow\varphi$,
stosujemy twierdzenie o zbieżności monotonicznej najpierw
w całce po $x$, a następnie w całkach po $y$ i po $\omega$.
Otrzymujemy wskazaną tożsamość dla każdej $\varphi\geq0$.

W szczególności dla $\varphi=|h|$ i $B=\R$ mamy
$\int_\R G_{|h|}\,\dd\Law(Y)=\E|h(X)|<\infty$.
Zatem borelowski zbiór $N=\{G_{|h|}=\infty\}$ jest
$\Law(Y)$-zerowy. Poza $N$ funkcje $G_{h^+}$ i $G_{h^-}$
są skończone, a $g=G_{h^+}-G_{h^-}$. Przypisanie wartości
zero na $N$ zachowuje borelowską mierzalność.
Ponadto $|g|\leq G_{|h|}$. Korzystając z
\atlasref{ME-13::thm:pushforward-integral-13}{wzoru na całkę względem miary obrazowej},
otrzymujemy
\[
\E|g(Y)|=\int_\R |g(y)|\,\dd\Law(Y)(y)
\leq\E|h(X)|<\infty.
\]
Funkcja $g(Y)$ jest $\sigma(Y)$-mierzalna, gdyż $g$ jest
borelowska. Odejmując tożsamości dla $h^+$ i $h^-$,
otrzymujemy dla każdego borelowskiego $B$
\[
\int_{Y^{-1}(B)}h(X)\,\dd\Prob
=\int_B g(y)\,\dd\Law(Y)(y)
=\int_{Y^{-1}(B)}g(Y)\,\dd\Prob.
\]
Rodzina $\{Y^{-1}(B):B\in\mathcal B(\R)\}$ jest
$\sigma$-ciałem i z definicji jest równa $\sigma(Y)$.
Sprawdziliśmy zatem mierzalność, całkowalność oraz
tożsamość całek wymaganą w definicji warunkowej wartości
oczekiwanej; por. rozdział~\ref{ch:condexp} i
\atlasref{ME-18::def:conditional-exp-18}{definicję w teorii miary i całki}.
Jednoznaczność prawie na pewno
(\atlasref{ME-18::thm:CE-existence-18}{twierdzenie o warunkowej wartości oczekiwanej})
daje tezę.
\end{proof}
\begin{przyklad}[Obliczanie warunkowej średniej i wariancji]
W modelu jednostajnym na trójkącie $0<x<y<1$
mamy $k(y,x)=y^{-1}\chi_{(0,y)}(x)$ dla $0<y<1$.
Ponieważ $0<X<Y<1$ prawie na pewno, zarówno $X$, jak i $X^2$
są całkowalne.
Zatem
$\E[X\mid Y=y]=y^{-1}\int_0^y x\,\dd\lambda_1(x)=y/2$
oraz
$\E[X^2\mid Y=y]=y^{-1}\int_0^y x^2\,\dd\lambda_1(x)=y^2/3$.
Wariancja rozkładu warunkowego przy $0<y<1$ jest równa
$y^2/3-(y/2)^2=y^2/12$, na mocy tożsamości wariancji
jako różnicy drugiego momentu i kwadratu pierwszego momentu.
Ponieważ $\Prob(0<Y<1)=1$, z powyższego twierdzenia dla
$h(x)=x$ oraz $h(x)=x^2$ i wzoru na wariancję warunkową
(rozdział~\ref{ch:condexp}) otrzymujemy poniższe równości.
Mianowicie $\E[X\mid\sigma(Y)]=Y/2$ i
$\Var(X\mid\sigma(Y))=Y^2/12$ prawie na pewno.
To przykład, w którym pojedyncza obserwacja $Y$
zmienia nie tylko prognozę, lecz również pozostałą
niepewność co do $X$.
\end{przyklad}

\section{Przekształcanie zmiennych i suma niezależnych zmiennych}
Jeśli znamy rozkład $X$, to często chcemy poznać rozkład
$g(X)$, na przykład kwadratu, logarytmu lub odwrotności.
Dla borelowskiego $g:\R\to\R$ zmienna $Y=g\circ X$ jest
mierzalna, a jej rozkład jest miarą obrazową:
$\mu_Y(B)=\mu_X(g^{-1}(B))$ dla $B\in\mathcal B(\R)$.
Wynika to z równości $\{Y\in B\}=\{X\in g^{-1}(B)\}$.
Przy zmianie zmiennych gęstość musi być skorygowana o to,
jak $g$ rozciąga albo ściska małe przedziały: stąd pochodna
funkcji odwrotnej. Przy sumie niezależnych zmiennych na
ustaloną wartość $s$ składają się wszystkie pary $(x,s-x)$;
„sumowanie” ich wkładów prowadzi do splotu.

\begin{stwierdzenie}[Transformacja monotoniczna]
Niech $X$ ma gęstość $f_X$, a $g:I\to J$ będzie
dyfeomorfizmem klasy $C^1$ przedziałów otwartych
$I,J\subseteq\R$. Załóżmy, że $\Prob(X\in I)=1$.
Definiujemy $Y=g(X)$ na $\{X\in I\}$, a poza tym
zdarzeniem przyjmujemy $Y=0$. Jest to zmienna losowa,
gdyż przedłużenie $g$ przez zero poza $I$ jest borelowskie.
Wtedy $Y$ ma gęstość
$f_Y(y)=f_X(g^{-1}(y))\,|(g^{-1})'(y)|$ dla $y\in J$
i gęstość zero poza $J$.
\end{stwierdzenie}
\begin{proof}
Dla dowolnego borelowskiego zbioru $B\subseteq\R$
rozpoczynamy od definicji rozkładu zmiennej $Y$.
W poniższych wzorach $g^{-1}$ oznacza odwrotność
pierwotnego dyfeomorfizmu $g:I\to J$: 
\[
\Prob(Y\in B)=\Prob(X\in g^{-1}(B\cap J))
=\int_{g^{-1}(B\cap J)}f_X(x)\,\dd\lambda_1(x).
\]
Pierwsza równość zachodzi, ponieważ zdarzenia po obu
stronach mogą różnić się tylko na $\{X\notin I\}$, którego
prawdopodobieństwo jest zerowe. Druga wynika z definicji
gęstości i borelowskości przeciwobrazu w $I$.
Do dyfeomorfizmu $g^{-1}:J\to I$ oraz funkcji
$x\mapsto f_X(x)\chi_{g^{-1}(B\cap J)}(x)$ stosujemy
\atlasref{ME-13::thm:change-variables-13}{twierdzenie o zamianie zmiennych w całce Lebesgue'a}.
Otrzymujemy
\[
\Prob(Y\in B)=\int_{B\cap J}
 f_X(g^{-1}(y))\,|(g^{-1})'(y)|\,\dd\lambda_1(y).
\]
Wskazana funkcja jest borelowska na $J$ jako iloczyn
funkcji borelowskiej i ciągłej, a jej przedłużenie zerem
poza $J$ jest również borelowskie. Jest nieujemna, a
dla $B=\R$ powyższa równość daje całkę równą jeden.
Zatem jest gęstością rozkładu $Y$.
Wartość bezwzględna przy pochodnej jest konieczna,
ponieważ odwracanie orientacji przedziału nie może
prowadzić do ujemnej gęstości.
\end{proof}

\begin{stwierdzenie}[Skończenie wiele gałęzi]
Niech $X$ ma gęstość $f_X$ i niech $I_1,\ldots,I_r$,
$r\geq1$, będą parami rozłącznymi przedziałami otwartymi
spełniającymi $\Prob(X\in\bigcup_{j=1}^r I_j)=1$.
Niech $g$ będzie funkcją na $\bigcup_{j=1}^r I_j$, której
ograniczenie do każdego $I_j$ jest $C^1$-dyfeomorfizmem
na otwarty przedział $J_j$.
Przedłużamy $g$ przez zero poza sumą tych przedziałów.
Dla $x\in I_j$ pochodna $g'(x)$ oznacza pochodną
ograniczenia $g|_{I_j}$; nie zakładamy istnienia pochodnej
na dopełnieniu sumy przedziałów.
Tak określona funkcja jest borelowska, gdyż jest ciągła
na każdym $I_j$ i stała na borelowskim dopełnieniu ich sumy.
Wtedy $Y=g(X)$ ma gęstość
\[
f_Y(y)=\sum_{j:y\in J_j}
\frac{f_X((g|_{I_j})^{-1}(y))}
{|g'((g|_{I_j})^{-1}(y))|}
\quad\text{dla prawie wszystkich }y\in\R.
\]
\end{stwierdzenie}
\begin{proof}
Niech $B\subseteq\R$ będzie borelowski. Ponieważ
$\Prob(X\in\bigcup_j I_j)=1$ i przedziały $I_j$
są rozłączne, możemy rozłożyć zdarzenie $\{g(X)\in B\}$
według gałęzi:
\[
\Prob(g(X)\in B)=\sum_{j=1}^r
\int_{I_j\cap g^{-1}(B)} f_X(x)\,\dd\lambda_1(x).
\]
Dla każdego $j$ oznaczmy $g_j=g|_{I_j}$.
Stosujemy twierdzenie o zamianie zmiennych osobno do
całki po $I_j\cap g^{-1}(B)$, bez normalizowania
ograniczonej miary. Z reguły łańcuchowej
$g_j'(g_j^{-1}(y))(g_j^{-1})'(y)=1$ dla $y\in J_j$,
więc $g_j'$ nie znika i
$|(g_j^{-1})'(y)|=1/|g_j'(g_j^{-1}(y))|$.
Otrzymujemy
\[
\Prob(g(X)\in B)=\sum_{j=1}^r
\int_{B\cap J_j}
\frac{f_X((g|_{I_j})^{-1}(y))}
{|g'((g|_{I_j})^{-1}(y))|}\,\dd\lambda_1(y).
\]
Każdy składnik przedłużamy przez zero poza $J_j$;
w szczególności nie obliczamy odwrotności gałęzi poza jej
dziedziną. Suma skończona jest nieujemna i borelowska,
a przeniesienie jej pod całkę daje podany wzór.
Dla $B=\R$ całka z tej funkcji wynosi jeden.
Wszystkie gałęzie prowadzące do tej samej wartości $y$
muszą zostać uwzględnione, ponieważ równanie $g(x)=y$
może mieć po jednym rozwiązaniu w kilku różnych $I_j$.
\end{proof}

\begin{twierdzenie}[Splot]
Niech $X,Y$ będą niezależnymi rzeczywistymi zmiennymi losowymi
o gęstościach $f_X,f_Y$. Wtedy $S=X+Y$ ma gęstość
\begin{equation}\label{eq:convolution}
 f_S(s)=(f_X*f_Y)(s)
 :=\int_\R f_X(x)f_Y(s-x)\,\dd\lambda_1(x)
\end{equation}
dla $\lambda_1$-prawie każdego $s$, przy konwencji
$0\cdot\infty=0$ dla gęstości o wartościach rozszerzonych.
\end{twierdzenie}
\begin{proof}
Z niezależności rozkład pary $(X,Y)$ jest iloczynem
rozkładów brzegowych. Z dowodu kryterium iloczynowego
gęstości (poprzednia sekcja o rozkładach łącznych) miara
produktowa ma gęstość $f_{X,Y}(x,y)=f_X(x)f_Y(y)$:
całka z iloczynu definiuje miarę probabilistyczną
i na prostokątach ma te same wartości co
$\mu_X\otimes\mu_Y$. Wybieramy wszędzie skończone
borelowskie wersje obu gęstości.
Niech $B\subseteq\R$ będzie dowolnym zbiorem borelowskim.
Funkcja dodawania jest ciągła, więc $S=X+Y$ jest mierzalna
jako złożenie z wektorem losowym $(X,Y)$.
Z definicji rozkładu łącznego
\[
\Prob(S\in B)=\int_{\R^2}\chi_B(x+y)
 f_X(x)f_Y(y)\,\dd\lambda_2(x,y).
\]
Rozważmy liniową zmianę współrzędnych
$(x,y)\mapsto (x,s=x+y)$, której macierz
$\left(\begin{smallmatrix}1&0\\1&1\end{smallmatrix}\right)$
ma wyznacznik równy jeden. Jej odwrotność to
$(x,s)\mapsto(x,s-x)$, również o wyznaczniku jeden.
Stosując \atlasref{ME-13::thm:change-variables-13}{twierdzenie o zamianie zmiennych}
do tej odwrotności, otrzymujemy
\[
\Prob(S\in B)=\int_{\R^2}\chi_B(s)
 f_X(x)f_Y(s-x)\,\dd\lambda_2(x,s).
\]
Ponieważ funkcja podcałkowa jest nieujemna,
możemy zmienić kolejność całkowania:
\[
\Prob(S\in B)=\int_B\left(
\int_\R f_X(x)f_Y(s-x)\,\dd\lambda_1(x)
\right)\dd\lambda_1(s).
\]
Wewnętrzna całka jest borelowską funkcją zmiennej $s$
na mocy twierdzenia Tonellego, ponieważ $(x,s)\mapsto
f_X(x)f_Y(s-x)$ jest nieujemna i borelowska.
Dla $B=\R$ otrzymujemy, że całka z tej funkcji wynosi jeden.
Jest ona zatem skończona $\lambda_1$-prawie wszędzie.
Z definicji gęstości wynika \eqref{eq:convolution}.
W razie potrzeby na borelowskim zbiorze miary zero,
na którym splot jest nieskończony, przyjmujemy wartość zero.
Wybór wersji gęstości $f_X,f_Y$ nie wpływa na samą całkę
definiującą splot dla żadnego ustalonego $s$ (także gdy
jej wartość jest nieskończona). Istotnie, jeśli wersje
różnią się na borelowskich zbiorach zerowych $N_X,N_Y$,
to dla ustalonego $s$ funkcje podcałkowe mogą różnić się
tylko na $N_X\cup(s-N_Y)$. Zbiór ten ma miarę Lebesgue'a
zero na mocy niezmienniczości względem odbicia i przesunięcia
(\atlasref{ME-04::thm:invariance-lebesgue}{niezmienniczość miary Lebesgue'a}).
Całki są zatem równe. Ewentualna zamiana wartości nieskończonych
splotu na zero zmienia jedynie wersję jego gęstości.
\end{proof}

\begin{przyklad}[Suma dwóch jednostajnych zmiennych]
Jeśli $U,V$ są niezależne i jednostajne na $(0,1)$,
to $f_U=f_V=\chi_{(0,1)}$. Z \eqref{eq:convolution}
dla $S=U+V$ otrzymujemy
\[
f_S(s)=\int_\R\chi_{(0,1)}(u)\chi_{(0,1)}(s-u)\,\dd\lambda_1(u)
=\lambda_1((0,1)\cap(s-1,s)).
\]
Długość przecięcia wynosi
$\max\{0,\min(1,s)-\max(0,s-1)\}$.
Po rozpatrzeniu zakresów $s\leq0$, $0<s\leq1$,
$1<s<2$ i $s\geq2$ dostajemy
\[
f_S(s)=\begin{cases}
s,&0<s\leq1,\\
2-s,&1<s<2,\\
0,&\text{w pozostałych przypadkach}.
\end{cases}
\]
Gęstość jest trójkątna, a nie stała na $(0,2)$.
Jej wartości odpowiadają długości zbioru dopuszczalnych
wartości pierwszej współrzędnej $u$, gdy suma wynosi $s$.
Całka z gęstości wynosi
$\int_0^1 s\,\dd\lambda_1(s)+\int_1^2(2-s)\,\dd\lambda_1(s)=1$.
\end{przyklad}

\begin{przyklad}[Bez niezależności wzór na splot może zawieść]
Niech $U\sim\Unif(0,1)$ oraz $V=1-U$.
Przekształcenie $u\mapsto1-u$ jest $C^1$-dyfeomorfizmem
$(0,1)$ na $(0,1)$, którego odwrotność ma pochodną $-1$.
Z twierdzenia o transformacji monotonicznej wynika
$V\sim\Unif(0,1)$. Tymczasem $U+V=1$ dla każdego wyniku. Rozkład sumy jest więc skupiony
w jednym punkcie i nie ma gęstości względem miary
Lebesgue'a, mimo że obie zmienne mają gęstości.
Nie jest to sprzeczne z twierdzeniem o splocie,
ponieważ $U$ i $V$ nie są niezależne:
$\Prob(U<1/2)=\Prob(V<1/2)=1/2$, ale
$\Prob(U<1/2,V<1/2)=0\neq1/4$.
\end{przyklad}


